\documentclass{amsart}
% For dealing with references we use the comment environment
\usepackage{verbatim}
\newenvironment{reference}{\comment}{\endcomment}
%\newenvironment{reference}{}{}
\newenvironment{slogan}{\comment}{\endcomment}
\newenvironment{history}{\comment}{\endcomment}

% For commutative diagrams we use Xy-pic
\usepackage[all]{xy}

% We use 2cell for 2-commutative diagrams.
\xyoption{2cell}
\UseAllTwocells

% We use multicol for the list of chapters between chapters
\usepackage{multicol}

% This is generally recommended for better output
\usepackage{lmodern}
\usepackage[T1]{fontenc}

% For cross-file-references

% Package for hypertext links:
\usepackage{hyperref}

% For any local file, say "hello.tex" you want to link to please

% Theorem environments.
%
\theoremstyle{plain}
\newtheorem{theorem}[subsection]{Theorem}
\newtheorem{proposition}[subsection]{Proposition}
\newtheorem{lemma}[subsection]{Lemma}

\theoremstyle{definition}
\newtheorem{definition}[subsection]{Definition}
\newtheorem{example}[subsection]{Example}
\newtheorem{exercise}[subsection]{Exercise}
\newtheorem{situation}[subsection]{Situation}

\theoremstyle{remark}
\newtheorem{remark}[subsection]{Remark}
\newtheorem{remarks}[subsection]{Remarks}

\numberwithin{equation}{subsection}

% Macros
%
\def\lim{\mathop{\mathrm{lim}}\nolimits}
\def\colim{\mathop{\mathrm{colim}}\nolimits}
\def\Spec{\mathop{\mathrm{Spec}}}
\def\Hom{\mathop{\mathrm{Hom}}\nolimits}
\def\Ext{\mathop{\mathrm{Ext}}\nolimits}
\def\SheafHom{\mathop{\mathcal{H}\!\mathit{om}}\nolimits}
\def\SheafExt{\mathop{\mathcal{E}\!\mathit{xt}}\nolimits}
\def\Sch{\mathit{Sch}}
\def\Mor{\mathop{\mathrm{Mor}}\nolimits}
\def\Ob{\mathop{\mathrm{Ob}}\nolimits}
\def\Sh{\mathop{\mathit{Sh}}\nolimits}
\def\NL{\mathop{N\!L}\nolimits}
\def\CH{\mathop{\mathrm{CH}}\nolimits}
\def\proetale{{pro\text{-}\acute{e}tale}}
\def\etale{{\acute{e}tale}}
\def\QCoh{\mathit{QCoh}}
\def\Ker{\mathop{\mathrm{Ker}}}
\def\Im{\mathop{\mathrm{Im}}}
\def\Coker{\mathop{\mathrm{Coker}}}
\def\Coim{\mathop{\mathrm{Coim}}}

% Boxtimes
%
\DeclareMathSymbol{\boxtimes}{\mathbin}{AMSa}{"02}

%
% Macros for moduli stacks/spaces
%
\def\QCohstack{\mathcal{QC}\!\mathit{oh}}
\def\Cohstack{\mathcal{C}\!\mathit{oh}}
\def\Spacesstack{\mathcal{S}\!\mathit{paces}}
\def\Quotfunctor{\mathrm{Quot}}
\def\Hilbfunctor{\mathrm{Hilb}}
\def\Curvesstack{\mathcal{C}\!\mathit{urves}}
\def\Polarizedstack{\mathcal{P}\!\mathit{olarized}}
\def\Complexesstack{\mathcal{C}\!\mathit{omplexes}}
% \Pic is the operator that assigns to X its picard group, usage \Pic(X)
% \Picardstack_{X/B} denotes the Picard stack of X over B
% \Picardfunctor_{X/B} denotes the Picard functor of X over B
\def\Pic{\mathop{\mathrm{Pic}}\nolimits}
\def\Picardstack{\mathcal{P}\!\mathit{ic}}
\def\Picardfunctor{\mathrm{Pic}}
\def\Deformationcategory{\mathcal{D}\!\mathit{ef}}

\setlength{\emergencystretch}{3em}
\begin{document}
\title[Stacks Project, Tag 07PV]{349 --- Essentially finite-type algebras over G-rings remain G-rings}
\maketitle
\noindent Copyright (C) 2005--2025 Johan de Jong. Stacks Project authors. Source: \href{https://stacks.math.columbia.edu/tag/07PV}{Tag 07PV}.

\noindent Editorial difficulty note (not author text). Difficulty H2: The proof reduces arbitrary finite-type algebras to one polynomial variable and compares local completions by a two-step square. It reduces formal fibres to a regular complete local domain and separates positive-characteristic prime cases, retaining the geometric-regularity obligation..

\noindent Editorial provenance note (not author text). History: excerpt from revision \texttt{a0444\allowbreak 6e57e\allowbreak c1fbc\allowbreak 252a8\allowbreak 71afc\allowbreak ec775\allowbreak 2fb28\allowbreak 07b14}, more-algebra.tex, lines 12551--12640. Reference presentation was adapted; original-author context and core support blocks were retained or added. The mathematical wording is unchanged. Licensed under GNU FDL 1.2 or later; no invariant sections or cover texts. See ../licenses/COPYING. Context blocks reproduce author definitions and supporting results; they are not additional counted problems.

\begin{definition}
\label{definition-G-ring}
A ring $R$ is called a {\it G-ring} if $R$ is Noetherian and for every
prime $\mathfrak p$ of $R$ the ring map
$R_\mathfrak p \to (R_\mathfrak p)^\wedge$ is regular.
\end{definition}

\begin{proposition}
\label{proposition-finite-type-over-G-ring}
Let $R$ be a G-ring. If $R \to S$ is essentially of finite type
then $S$ is a G-ring.
\end{proposition}

\begin{proof}
Since being a G-ring is a property of the local rings it is clear
that a localization of a G-ring is a G-ring. Conversely, if every
localization at a prime is a G-ring, then the ring is a G-ring.
Thus it suffices to show that $S_\mathfrak q$ is a G-ring for every
finite type $R$-algebra $S$ and every prime $\mathfrak q$ of $S$.
Writing $S$ as a quotient of $R[x_1, \ldots, x_n]$ we see from
Lemma \href{https://stacks.math.columbia.edu/tag/07PP}{lemma G ring goes up quasi finite (Tag 07PP)} that it suffices to prove
that $R[x_1, \ldots, x_n]$ is a G-ring. By induction on $n$ it
suffices to prove that $R[x]$ is a G-ring. Let $\mathfrak q \subset R[x]$
be a maximal ideal. By Lemma \href{https://stacks.math.columbia.edu/tag/07PT}{lemma check G ring maximal ideals (Tag 07PT)}
it suffices to show that
$$
R[x]_\mathfrak q \longrightarrow R[x]_\mathfrak q^\wedge
$$
is regular. If $\mathfrak q$ lies over $\mathfrak p \subset R$, then
we may replace $R$ by $R_\mathfrak p$. Hence we may assume that $R$
is a Noetherian local G-ring with maximal ideal $\mathfrak m$ and
that $\mathfrak q \subset R[x]$ lies over $\mathfrak m$. Note that
there is a unique prime $\mathfrak q' \subset R^\wedge[x]$
lying over $\mathfrak q$. Consider the diagram
$$
\xymatrix{
R[x]_\mathfrak q^\wedge \ar[r] &
(R^\wedge[x]_{\mathfrak q'})^\wedge \\
R[x]_\mathfrak q \ar[r] \ar[u] & R^\wedge[x]_{\mathfrak q'} \ar[u]
}
$$
Since $R$ is a G-ring the lower horizontal arrow is regular
(as a localization of a base change of the regular ring map
$R \to R^\wedge$). Suppose we can prove the right vertical arrow
is regular. Then it follows that the composition
$R[x]_\mathfrak q \to (R^\wedge[x]_{\mathfrak q'})^\wedge$
is regular, and hence the left vertical arrow is regular by
Lemma \href{https://stacks.math.columbia.edu/tag/07NT}{lemma regular permanence (Tag 07NT)}.
Hence we see that we may assume $R$ is a Noetherian complete
local ring and $\mathfrak q$ a prime lying over the maximal
ideal of $R$.

\medskip\noindent
Let $R$ be a Noetherian complete local ring and let $\mathfrak q \subset R[x]$
be a maximal ideal lying over the maximal ideal of $R$. Let
$\mathfrak r \subset \mathfrak q$ be a prime ideal. We want to show that
$R[x]_\mathfrak q^\wedge \otimes_{R[x]} \kappa(\mathfrak r)$ is
a geometrically regular algebra over $\kappa(\mathfrak r)$.
Set $\mathfrak p = R \cap \mathfrak r$. Then we can replace $R$
by $R/\mathfrak p$ and $\mathfrak q$ and $\mathfrak r$ by their
images in $R/\mathfrak p[x]$, see
Lemma \href{https://stacks.math.columbia.edu/tag/07PN}{lemma check G ring easy (Tag 07PN)}.
Hence we may assume that $R$ is a domain and that $\mathfrak r \cap R = (0)$.

\medskip\noindent
By  Algebra, Lemma
\href{https://stacks.math.columbia.edu/tag/032D}{lemma complete local Noetherian domain finite over regular (Tag 032D)}
we can find $R_0 \subset R$ which is regular and such that $R$ is
finite over $R_0$. Applying Lemma \href{https://stacks.math.columbia.edu/tag/07PP}{lemma G ring goes up quasi finite (Tag 07PP)}
we see that it suffices to prove
$R[x]_\mathfrak q^\wedge \otimes_{R[x]} \kappa(\mathfrak r)$
is geometrically regular over $\kappa(\mathfrak r)$ when, in addition to the
above, $R$ is a regular complete local ring.

\medskip\noindent
Now $R$ is a regular complete local ring, we have
$\mathfrak r \subset \mathfrak q \subset R[x]$, we have
$(0) = R \cap \mathfrak r$ and $\mathfrak q$ is a maximal ideal
lying over the maximal ideal of $R$. Since $R$ is regular the
ring $R[x]$ is regular (Algebra, Lemma \href{https://stacks.math.columbia.edu/tag/07NF}{lemma regular goes up (Tag 07NF)}).
Hence the localization $R[x]_\mathfrak q$ is regular.
Hence the completions $R[x]_\mathfrak q^\wedge$ are regular, see
Lemma \href{https://stacks.math.columbia.edu/tag/07NY}{lemma completion regular (Tag 07NY)}.
Hence the fibre $R[x]_{\mathfrak q}^\wedge \otimes_{R[x]} \kappa(\mathfrak r)$
is, as a localization of $R[x]_\mathfrak q^\wedge$, also regular.
Thus we are done if the characteristic of the fraction field of $R$ is $0$.

\medskip\noindent
If the characteristic of $R$ is positive, then $R = k[[x_1, \ldots, x_n]]$.
In this case we split the argument in two subcases:
\begin{enumerate}
\item The case $\mathfrak r = (0)$. The result is a direct consequence
of Lemma \href{https://stacks.math.columbia.edu/tag/07PR}{lemma helper G ring (Tag 07PR)}.
\item The case $\mathfrak r \not = (0)$. This is
Lemma \href{https://stacks.math.columbia.edu/tag/07PU}{lemma another helper G ring (Tag 07PU)}.
\end{enumerate}
\end{proof}
\end{document}
